\section{Related Work}
\label{sec:related}

\subsection{Model-based and traditional machine-learning methods}

Physics-based prognostics build degradation models from electrochemical
dynamics, estimated with filters such as extended Kalman filters
\cite{ref_ekf_yan}, unscented particle filters \cite{ref_upf_cong},
and particle filters with electrochemical or empirical state models
\cite{ref_pf_echem,ref_pf_rbf,ref_pf_lui}. These methods encode
mechanistic knowledge but require accurate parameter identification
and struggle with regeneration-rich or heterogeneous degradation.
Traditional machine learning provides a lighter alternative:
support vector machines \cite{ref_svm_klass,ref_svm_patil}, support
vector regression \cite{ref_svr_wang}, and relevance vector machines
\cite{ref_rvm_guo} regress health indicators directly from capacity
or voltage features, yet their fixed kernels cannot easily extrapolate
over the strongly nonlinear EOL-acceleration phase.

\subsection{Deep sequence models for battery prognostics}

Recurrent models were among the first deep approaches to capacity sequences:
LSTM-based predictors \cite{ref_lstm_rul_zhang,ref_lstm_elman_li},
GRU variants \cite{ref_gru_ding,ref_gru_rnn}, and recurrent networks
for RUL \cite{ref_rnn_catelani} capture temporal dynamics but suffer
from long-term dependency limitations. Convolutional hybrids combine
local feature extraction with recurrence, including CNN-BiLSTM
\cite{ref_cnn_bilstm}, TCN-LSTM \cite{ref_tcn_lstm}, and attention
mechanisms \cite{ref_lstm_transformer}. Transformer-based models
\cite{ref_vaswani} brought global attention to battery prognostics:
direct Transformer RUL predictors \cite{ref_trans_rul}, hybrid
CEEMDAN-Transformer-DNN models \cite{ref_ceemdan_trans}, and
PatchFormer \cite{ref_patchformer}, which combines dual patch-wise
attention with feature-wise attention. General time-series
transformers evaluated on battery tasks include Autoformer
\cite{ref_autoformer}, FEDformer \cite{ref_fedformer}, iTransformer
\cite{ref_itransformer}, PatchTST \cite{ref_patchtst}, TimesNet
\cite{ref_timesnet}, and TimeMixer \cite{ref_timemixer}. All
attention-based methods share the quadratic-complexity bottleneck
that constrains embedded deployment. Comprehensive surveys of this
line are given in \cite{ref_lipu_review,ref_he_review}.

\subsection{State space models}

Mamba \cite{ref_mamba} is a selective state space model with
linear-time inference, extensively surveyed in \cite{ref_mamba_survey}
and evaluated for general time-series forecasting in
\cite{ref_samamba_tsf}. For batteries, MambaLithium \cite{ref_mambalithium}
applies a pure Mamba network to SOH/RUL/SOC estimation, and RUL-Mamba
\cite{ref_rulmamba} combines a Mamba encoder with a feature attention
network and a gated residual decoder, reporting state-of-the-art
results on NASA, Oxford, and TJU datasets. While these works emphasize
efficient training and fast inference, their efficiency claims are
based on GPU measurements; none of them validates execution on
embedded hardware, which is the target of the present work.

\subsection{Grey system models}

Grey system theory models degradation using accumulated sequences and
fractional-order operators. Recent work includes grey models with
capacity-regeneration consideration \cite{ref_grey_hybrid}, adaptive
weighted fractional-order reverse accumulation for SOH
\cite{ref_grey_frac}, fractional-order heuristic models for PEMFC
prognostics \cite{ref_grey_pemfc}, self-adaptive grey neural networks
for real-time PEMFC degradation \cite{ref_grey_nn}, grey Kalman
filtering for SOC estimation \cite{ref_grey_kalman}, and grey
iterative modeling for battery health management \cite{ref_grey_iter}.
These methods are lightweight and interpretable but rely on
hand-crafted model forms that must be re-specified for each battery
chemistry and operating regime.

\subsection{Physics-informed battery modeling}

Physics-informed neural networks (PINNs) embed degradation physics
into the learning objective. Wang et al.\ \cite{ref_wang_pinn}
proposed a physics-informed network for stable degradation modeling,
and PiDDM \cite{ref_piddm} incorporates Arrhenius degradation kinetics
into the training objective, reporting improved extrapolation at the
end of life. Our systematic injection-route study (Section
\ref{sec:ablation}) revisits this paradigm for a sequence model that
already ingests the capacity trajectory: objective regularization,
input feature concatenation, and structural heads all fail---the
per-window z-score target space is structurally incompatible with
monotone predictors, and the features are redundant with capacity
history. The mechanism that survives is a degradation-rate head
driven by measured internal resistance, trained in absolute capacity
space (Section~\ref{sec:physics}).

\subsection{Linear attention and recurrent mixing}

Linear attention replaces the softmax attention map with a recurrent
matrix state, achieving $O(N)$ complexity. Gated DeltaNet
\cite{ref_gdn} combined the delta rule with learned gating, and
Gated DeltaNet-2 (GDN-2) \cite{ref_gdn2} decoupled channel-wise erase
and write gates, improving long-context retrieval in language
modeling. To our knowledge, the present work is the first to report
floating-point-accurate MCU deployment of a linear-attention battery model.

Stepping back from individual architectures, the surveyed literature
answers one question---which model is most accurate on public cycling
data---while leaving the questions that decide operational adoption
largely open. First, accuracy has been demonstrated only on
server-class hardware: none of the battery-specific models above
reports a deployment path onto the microcontrollers that production
BMS firmware actually runs on, so operators must either pay for
cloud connectivity or accept simpler on-device estimators. Second,
forecasts are evaluated as point trajectories, yet the operational
consumer of a forecast is a maintenance schedule, which requires an
explicit statement of risk; few works report calibrated intervals
suitable for scheduling. Third, evaluations are confined to clean
laboratory data: the end-of-life acceleration tail and the
corrupted-sensor regimes---where field decisions concentrate---are
rarely tested, so reported accuracy says little about how a
predictor will behave in deployment. DeltaCycle is designed around
these three operational gaps rather than around a single benchmark
number.
